\documentclass[11pt,a4paper]{article}
\input{preamble}

\pagestyle{fancy}
\fancyhf{}
\fancyhead[L]{\small\color{lightbrown}\emoji{robot} RL 틱택토 퀵가이드}
\fancyhead[R]{\small\color{lightbrown}v1.0.0.0}
\fancyfoot[C]{\thepage}
\renewcommand{\headrulewidth}{0.4pt}
\setlength{\headheight}{14pt}

\begin{document}

% === 타이틀 페이지 ===
\begin{titlepage}
\centering
\vspace*{3cm}
{\fontsize{48pt}{50pt}\selectfont \emoji{rocket}}
\vspace{1cm}

{\Huge\bfseries\color{mainblue} RL 틱택토}
\vspace{0.3cm}

{\LARGE\color{darkbrown} 퀵가이드}
\vspace{1.5cm}

{\large 시뮬레이터 옆에 띄워두고 참조하세요!}
\vspace{3cm}

\begin{tabular}{rl}
\textbf{문서 버전} & v1.0.0.0 \\
\textbf{시뮬레이터} & rl-tictactoe-v3.html (v3) \\
\textbf{작성일} & 2026-02-09 \\
\end{tabular}
\vfill
{\small\color{lightbrown} 청주교육대학교 — 컴퓨팅사고와 AI교육}
\end{titlepage}

\tableofcontents
\newpage

% ============================================================
\section{\emoji{books} 이 문서 사용법}
% ============================================================

\begin{center}
\renewcommand{\arraystretch}{1.4}
\begin{tabular}{|l|l|}
\hline
\rowcolor{tableheader}
\textbf{상황} & \textbf{참조 섹션} \\
\hline
처음 시작할 때 & \emoji{high-voltage} 초속성 실습 \\
\hline
AI를 학습시키고 싶을 때 & \emoji{rocket} 빠른 시작 \\
\hline
화면이 뭔지 모르겠을 때 & \emoji{desktop-computer} 화면 구성 \\
\hline
파라미터를 바꿔보고 싶을 때 & \emoji{gear} 파라미터 설명 \\
\hline
뭘 관찰해야 할지 모를 때 & \emoji{magnifying-glass-tilted-left} 관찰 포인트 \\
\hline
문제가 생겼을 때 & \emoji{sos} 트러블슈팅 \\
\hline
용어가 이해 안 될 때 & \emoji{key} 핵심 용어 \\
\hline
\end{tabular}
\end{center}

% ============================================================
\section{\emoji{high-voltage} 1분 초속성 실습}
% ============================================================

\begin{center}
\renewcommand{\arraystretch}{1.4}
\begin{tabular}{|c|p{7cm}|p{5cm}|}
\hline
\rowcolor{tableheader}
\textbf{단계} & \textbf{행동} & \textbf{확인} \\
\hline
1 & \textbf{[\emoji{play-button} 학습 시작]} 버튼을 클릭하세요 & 그래프가 움직이기 시작하나요? \\
\hline
2 & 30초 기다리세요 (속도 \textbf{1ms} 상태) & 에피소드 수가 빠르게 올라가나요? \\
\hline
3 & \textbf{[\emoji{bust-in-silhouette} 사람 대전]} 버튼을 클릭하세요 & $3\times3$ 보드가 크게 나타나나요? \\
\hline
4 & 빈 칸을 클릭해서 돌을 놓으세요 & AI가 바로 응수하나요? \emoji{party-popper} \\
\hline
\end{tabular}
\end{center}

\begin{tipbox}[\emoji{light-bulb} 팁]
AI를 이긴다면 아직 학습이 부족한 것입니다! 더 학습시켜 보세요.
\end{tipbox}

% ============================================================
\section{\emoji{rocket} 빠른 시작 (3단계)}
% ============================================================

\subsection{1단계: AI 학습시키기}

파일(\texttt{rl-tictactoe-v3.html})을 브라우저에서 열고, 좌측 패널에서 \textbf{\emoji{clipboard} 학습 방식}을 확인합니다 (기본: \emoji{counterclockwise-arrows-button} 자기대국). \textbf{[\emoji{play-button} 학습 시작]} 클릭하면 우측 그래프가 실시간으로 변합니다.

\begin{infobox}[\emoji{bullseye} 목표]
\textbf{5,000판} 이상이면 효과가 보이기 시작하고, \textbf{20,000판} 이상이면 최적에 가까워집니다.
\end{infobox}

\subsection{2단계: AI와 대국하기}

\textbf{[\emoji{bust-in-silhouette} 사람 대전]} 클릭하면 대국 화면으로 전환됩니다. 보드의 빈 칸을 클릭하여 돌을 놓으면 AI가 즉시 응수합니다. 오른쪽 \textbf{\emoji{brain} AI 판단 근거} 히트맵을 확인하세요.

\subsection{3단계: 실험하기}

대시보드로 돌아와서 학습을 정지하고, 좌측 \textbf{\emoji{gear} 파라미터}에서 학습률($\alpha$)을 바꿔봅니다. 다시 학습을 시작하면 그래프 변화를 관찰할 수 있습니다.

% ============================================================
\section{\emoji{desktop-computer} 화면 구성}
% ============================================================

\subsection{대시보드 (메인 화면)}

\begin{center}
\renewcommand{\arraystretch}{1.4}
\begin{tabular}{|l|p{9cm}|}
\hline
\rowcolor{tableheader}
\textbf{영역} & \textbf{구성 요소} \\
\hline
좌측 패널 (조작) & 미니보드, 속도 슬라이더, 학습 방식, 파라미터($\alpha$/$\gamma$/$\varepsilon$), 액션 버튼, 저장/불러오기, 초기화 \\
\hline
우측 대시보드 (관찰) & 학습 단계 진행바(\emoji{seedling}$\to$\emoji{magnifying-glass-tilted-left}$\to$\emoji{books}$\to$\emoji{bullseye}$\to$\emoji{trophy}), 지표 카드 6개, 그래프 4개, 인사이트 패널 \\
\hline
\end{tabular}
\end{center}

\subsection{사람 대전}

\begin{center}
\renewcommand{\arraystretch}{1.4}
\begin{tabular}{|l|p{9cm}|}
\hline
\rowcolor{tableheader}
\textbf{영역} & \textbf{구성 요소} \\
\hline
좌측 & $3\times3$ 게임 보드, 차례 표시, 새 게임/선후공 전환, 승무패 기록 \\
\hline
우측 & 학습 정보 패널, Q-값 히트맵 (AI 판단 근거) \\
\hline
\end{tabular}
\end{center}

% ============================================================
\section{\emoji{gear} 파라미터 설명}
% ============================================================

\begin{center}
\renewcommand{\arraystretch}{1.4}
\begin{tabular}{|l|c|p{6.5cm}|}
\hline
\rowcolor{tableheader}
\textbf{파라미터} & \textbf{기본값} & \textbf{쉬운 설명} \\
\hline
학습률 $\alpha$ & 0.30 & 한 번의 경험에서 얼마나 배울지. 높으면 빠르지만 불안정 \\
\hline
할인율 $\gamma$ & 0.90 & 미래 보상을 얼마나 중요하게 볼지. 높으면 장기적 전략 \\
\hline
$\varepsilon$ 감소율 & 0.9995 & 탐험을 줄이는 속도. 1에 가까울수록 천천히 줄임 \\
\hline
속도 & 1ms & 학습 속도. 1ms가 가장 빠름 (500판/틱) \\
\hline
\end{tabular}
\end{center}

\subsection{파라미터 실험 가이드}

\begin{center}
\renewcommand{\arraystretch}{1.4}
\begin{tabular}{|p{3.5cm}|p{4cm}|p{4.5cm}|}
\hline
\rowcolor{tableheader}
\textbf{실험} & \textbf{설정} & \textbf{관찰 포인트} \\
\hline
빠른 학습 vs 느린 학습 & $\alpha$를 0.9 $\to$ 0.01로 변경 & 승률 그래프의 수렴 속도 차이 \\
\hline
탐험 많이 vs 적게 & $\varepsilon$ 감소율을 0.99 $\to$ 0.9999로 & 탐험률 그래프, 학습 후반 성능 \\
\hline
단기 vs 장기 전략 & $\gamma$를 0.1 $\to$ 0.99로 변경 & AI 플레이 스타일 변화 \\
\hline
\end{tabular}
\end{center}

% ============================================================
\section{\emoji{magnifying-glass-tilted-left} 이런 장면을 찾아보세요!}
% ============================================================

\begin{center}
\renewcommand{\arraystretch}{1.4}
\begin{tabular}{|l|p{4.5cm}|p{4.5cm}|}
\hline
\rowcolor{tableheader}
\textbf{관찰 대상} & \textbf{찾아볼 장면} & \textbf{의미} \\
\hline
X 승률 그래프 & 처음 50\% 근처 $\to$ 점차 상승 & AI가 승리 전략을 학습 \\
\hline
무승부율 그래프 & 후반에 점점 증가 & 양쪽 AI 모두 실력 상승 \\
\hline
$\varepsilon$ 그래프 & 100\%에서 점차 하강 & 탐험$\to$활용으로 전환 \\
\hline
Q-값 히트맵 & 중앙(4번)이 가장 높은 값 & ``중앙이 가장 유리하다'' \\
\hline
첫수 선호도 & \emoji{bullseye} 중앙 > 모서리 > 변 & 틱택토의 최적 전략과 일치! \\
\hline
\end{tabular}
\end{center}

% ============================================================
\section{\emoji{bullseye} 실습 활동}
% ============================================================

\subsection{활동 1: AI의 성장 관찰}

\begin{center}
\renewcommand{\arraystretch}{1.4}
\begin{tabular}{|c|c|c|p{5cm}|}
\hline
\rowcolor{tableheader}
\textbf{학습량} & \textbf{AI 강도} & \textbf{대국 결과} & \textbf{관찰 내용} \\
\hline
1,000판 & \hspace{2cm} & \hspace{2cm} & \\
\hline
10,000판 & & & \\
\hline
50,000판 & & & \\
\hline
\end{tabular}
\end{center}

\subsection{활동 2: 파라미터의 영향}

\begin{center}
\renewcommand{\arraystretch}{1.4}
\begin{tabular}{|c|c|c|p{5cm}|}
\hline
\rowcolor{tableheader}
\textbf{학습률 $\alpha$} & \textbf{10,000판 후 강도} & \textbf{그래프 모양} & \textbf{관찰 내용} \\
\hline
0.01 & \hspace{2cm} & \hspace{2cm} & \\
\hline
0.50 & & & \\
\hline
0.99 & & & \\
\hline
\end{tabular}
\end{center}

% ============================================================
\section{\emoji{sos} 도와주세요!}
% ============================================================

\begin{center}
\renewcommand{\arraystretch}{1.4}
\small
\begin{tabular}{|p{3.5cm}|p{3cm}|p{5cm}|}
\hline
\rowcolor{tableheader}
\textbf{증상} & \textbf{원인} & \textbf{해결책} \\
\hline
화면이 안 뜸 & 인터넷 필요 (CDN) & 인터넷 연결 확인 후 새로고침 \\
\hline
그래프가 안 움직임 & 정지 상태 & [\emoji{play-button} 학습 시작] 다시 클릭 \\
\hline
AI가 너무 쉬움 & 학습 부족 & 속도 1ms로 20,000판 이상 학습 \\
\hline
AI를 이길 수 없음 & AI 최적 학습 완료 & 학습 초기화(\emoji{wastebasket}) 후 재학습 \\
\hline
데이터 사라짐 & 브라우저 초기화 & \emoji{floppy-disk} 저장 또는 \emoji{outbox-tray} 파일 내보내기 \\
\hline
그래프 요동침 & $\alpha$ 너무 높음 & $\alpha$를 0.1 이하로 낮추기 \\
\hline
글씨 안 보임 (프로젝터) & 다크 테마 & [\emoji{sun} 밝게] 버튼으로 전환 \\
\hline
\end{tabular}
\end{center}

% ============================================================
\section{\emoji{key} 핵심 용어}
% ============================================================

\begin{center}
\renewcommand{\arraystretch}{1.4}
\small
\begin{tabular}{|l|l|p{6cm}|}
\hline
\rowcolor{tableheader}
\textbf{용어} & \textbf{학술 용어} & \textbf{쉬운 설명} \\
\hline
Q-값 & Q-value & ``이 칸에 두면 얼마나 좋을까?'' 점수 \\
\hline
탐험률($\varepsilon$) & Epsilon & ``모험할까, 잘하는 걸 할까?'' 비율 \\
\hline
학습률($\alpha$) & Learning Rate & ``한 번 경험에서 얼마나 배울까?'' \\
\hline
할인율($\gamma$) & Discount Factor & ``나중 보상도 소중하게 여길까?'' \\
\hline
보상 & Reward & 이기면 +점수, 지면 $-$점수 \\
\hline
에피소드 & Episode & 한 판의 게임 (시작$\sim$종료) \\
\hline
Q-테이블 & Q-table & AI의 ``경험 수첩'' \\
\hline
정규화 & Canonicalization & 돌려도 같은 모양이면 같은 걸로 취급 \\
\hline
\end{tabular}
\end{center}

% ============================================================
\section{\emoji{thought-balloon} 생각해보기}
% ============================================================

\begin{infobox}[\emoji{red-question-mark} 질문 1]
틱택토에서 양쪽 모두 최선을 다하면 결과는 항상 무승부입니다. 그렇다면 ``강화학습이 성공했다''는 것은 무엇을 의미할까요?
\end{infobox}

\begin{infobox}[\emoji{red-question-mark} 질문 2]
AI가 학습 초반에 탐험(랜덤)을 많이 하는 이유는 무엇일까요? 처음부터 잘하는 걸 하면 안 될까요?
\end{infobox}

\begin{infobox}[\emoji{red-question-mark} 질문 3]
이 시뮬레이터의 AI는 Q-테이블에 경험을 기록합니다. 바둑이나 체스처럼 경우의 수가 훨씬 많은 게임에서도 같은 방법이 통할까요?
\end{infobox}

% ============================================================
\newpage
\section*{참고자료}
% ============================================================

\begin{center}
\renewcommand{\arraystretch}{1.4}
\begin{tabular}{|c|l|l|l|}
\hline
\rowcolor{tableheader}
\textbf{\#} & \textbf{자료명} & \textbf{유형} & \textbf{비고} \\
\hline
1 & RL틱택토\_개발일지\_v1.0.0.0 & 프로젝트 문서 & 시스템 설계, 기술 상세 \\
\hline
2 & 코딩관리지침\_v3.0.0.0 & 프로젝트 문서 & §6 퀵가이드 작성법 \\
\hline
3 & 버전관리지침\_v1.3.0 & 프로젝트 문서 & 버전 번호 체계 \\
\hline
\end{tabular}
\end{center}

\section*{생성/수정 이력}

\begin{center}
\renewcommand{\arraystretch}{1.3}
\begin{tabular}{|l|l|l|l|p{6cm}|l|}
\hline
\rowcolor{tableheader}
\textbf{버전} & \textbf{날짜} & \textbf{시간} & \textbf{변경 수준} & \textbf{변경 내용} & \textbf{작성자} \\
\hline
v1.0.0.0 & 2026-02-09 & 15:30 & 최초 & 초안 작성 — 필수 7섹션 + 선택 3섹션 & Claude \\
\hline
\end{tabular}
\end{center}

\vfill
\begin{center}
\small\color{lightbrown}
\emph{RL 틱택토 퀵가이드 — 시뮬레이터 옆에 띄워두고 참조하세요!}
\end{center}

\end{document}
